% © 2026 Ing. David Jaroš — CC BY-NC-ND 4.0
% Status: B-tier machine-verified note with selected claims human-reviewed and author-approved.
% UBT-REVIEW-PROFILE-BEGIN
% schema: ubt-review-profile/v1
% machine_verification: comprehensive_for_named_claims
% human_review: selected_claims
% editorial_approval: approved
% registry: ../../PROVENANCE_REVIEW.yaml
% UBT-REVIEW-PROFILE-END

% UBT-AI-PROVENANCE-BEGIN
% schema: ubt-ai-provenance/v1
% tier: B_machine_verified
% ai_assistance: disclosed
% human_review: machine-verification
% editorial_responsibility: Ing. David Jaroš
% policy: ../../AI_PROVENANCE.md
% notice: Machine-verified against named sources or verifiers; individual attestation is not claimed.
% UBT-AI-PROVENANCE-END

\documentclass[11pt,a4paper]{article}
\usepackage[margin=2.5cm]{geometry}
\usepackage{amsmath,amssymb,amsthm,mathtools}
\usepackage[most]{tcolorbox}
\usepackage[hidelinks]{hyperref}
\usepackage{ubtprovenance}
\UBTTier{B}
\UBTReviewProfile{comprehensive for named claims}{selected principal claims}{author approved}
\newtheorem{theorem}{Theorem}
\newtheorem{corollary}[theorem]{Corollary}
\newtheorem{proposition}[theorem]{Proposition}
\newtheorem{remark}[theorem]{Remark}
\newcommand{\Hess}{\mathcal H}
\newcommand{\Fcal}{\mathcal F}
\title{GAP-10D $\Theta$-Hessian Principal-Symbol Decision:\\
Fixed-Background Laplace Type, a Nonconic Resonance No-Go,\\
and the Exact Composite Remainder}
\author{Ing. David Jaro\v{s}}
\date{1 August 2026}
\begin{document}
\maketitle
\UBTProvenanceNotice

\begin{abstract}
We perform the cheapest decisive test proposed for the induced-gravity route,
but distinguish the quadratic Hessian of an action from the frozen
self-consistency symbol of the $D$-composite equation.  For the canonical
$\Theta$ kinetic term with metric, connection and nondegenerate internal
pairing held fixed, the raised quadratic Hessian has the scalar principal
symbol
\(
 g^{\mu\nu}k_\mu k_\nu\mathbf1
\).
After Euclidean continuation it is a minimal Laplace-type operator, so the
standard heat-kernel machinery applies to this fixed-background fluctuation
problem.  By contrast, the six-dimensional $q=1$ sector of the frozen
$D$-composite equation is a finite-scale singularity of the mixed-order full
symbol $I-A(s,\lambda)$.  Its singular surface is not conic under radial
rescaling of the covector, and the first-order principal symbol $-A$ is
generically rank nine.  Therefore that six-dimensional resonance cannot be
used directly as a Laplace bundle in a Gilkey calculation.  The remaining
first-principles task is now stated exactly: derive, gauge-fix and classify the
full composite $\Theta$-only Hessian including every chain-rule variation of
$E[\Theta]$, $g[\Theta]$ and $\Omega[\Theta]$.  No Einstein coefficient,
physical mode count or unconditional GR derivation is claimed here.  All
symbolic statements are checked by
\path{tools/verify_theta_hessian_principal_symbol.py}.
\end{abstract}

\section{Scope and the two operators that must not be conflated}
The working internal action contains a kinetic term of the schematic form
\begin{equation}
 S_\Theta=\frac12\int d^4x\sqrt{|g|}\,
 K_{AB}(D_\mu\Theta)^A(D^\mu\Theta)^B-\int d^4x\sqrt{|g|}\,V(\Theta),
 \label{eq:kinetic}
\end{equation}
where $K_{AB}$ is a nondegenerate real pairing on the selected fluctuation
space.  There are two inequivalent linear objects in the present repository:
\begin{enumerate}
 \item the Hessian $\Hess=\delta^2 S/\delta\Theta^2$ of an action;
 \item the frozen self-consistency symbol
 \((I-A(s,\lambda))F=\partial\theta'$ of the torsion-free
 $D$-composite defining equation.
\end{enumerate}
Only the first object is an input to a one-loop determinant.  A singularity of
the second object may identify a kinematic resonance, but it is not thereby a
quadratic action or a Laplace-type Hessian.

\section{Fixed-background Hessian}
Hold $g_{\mu\nu}$, the background connection and the internal pairing fixed,
and write $\Theta=\bar\Theta+\vartheta$.  Let
$D_\mu=\partial_\mu+C_\mu$ on the real fluctuation bundle.  Up to boundary
terms, the quadratic action has the form
\begin{equation}
 S^{(2)}_{\rm fixed}
 =\frac12\int d^4x\sqrt{|g|}\left[
 K_{AB}(D_\mu\vartheta)^A(D^\mu\vartheta)^B
 +\vartheta^A U_{AB}\vartheta^B\right],
 \label{eq:quadratic}
\end{equation}
where $U$ includes the potential Hessian and other derivative-free terms.

\begin{theorem}[Fixed-background scalar principal symbol]
\label{thm:fixed}
If $K$ is nondegenerate, the lower-index Hessian of
Eq.~\eqref{eq:quadratic} has
\begin{equation}
 \sigma_2(\Hess_{AB})(x,k)
 =K_{AB}\,g^{\mu\nu}(x)k_\mu k_\nu.
 \label{eq:lower-symbol}
\end{equation}
After raising one internal index with $K^{-1}$,
\begin{equation}
 \boxed{\sigma_2(\Hess^A{}_B)(x,k)
 =g^{\mu\nu}(x)k_\mu k_\nu\,\delta^A{}_B.}
 \label{eq:scalar-symbol}
\end{equation}
The background connection contributes only first- and zero-order terms, and
$U$ contributes only zero order.
\end{theorem}

\begin{proof}
In the highest-derivative part of Eq.~\eqref{eq:quadratic}, replace
$D_\mu$ by $\partial_\mu$.  Every insertion of $C_\mu$ lowers the number of
derivatives acting on $\vartheta$.  Integration by parts therefore gives
$-K_{AB}g^{\mu\nu}\partial_\mu\partial_\nu$ plus terms of order at most one.
The potential Hessian is derivative-free.  Replacing
$\partial_\mu\mapsto ik_\mu$ and raising the internal index proves
Eq.~\eqref{eq:scalar-symbol}.  The exact polynomial extraction is reproduced
by the verifier.
\end{proof}

\begin{corollary}[Conditional heat-kernel admissibility]
After a controlled Euclidean continuation, the fixed-background operator can
be written in minimal Laplace form
\begin{equation}
 \Hess_E=-\bigl(g_E^{\mu\nu}\nabla_\mu\nabla_\nu+\mathcal E\bigr).
\end{equation}
Consequently the standard local heat-kernel coefficients apply to this
specified fluctuation problem.  This proves the operator-class assumption
used in the companion induced-gravity calculation; it does not prove that the
same operator is the full composite gravitational Hessian.
\end{corollary}

\begin{proposition}[Metric-lock collapse of the same kinetic scalar]
\label{prop:lock}
If the metric is no longer fixed but is defined by the same covariant jet,
\[
 E_\mu=\mathcal N_0^{-1/2}D_\mu\Theta,
 \qquad
 \langle E_\mu,E_\nu\rangle_\sharp=g_{\mu\nu},
\]
then in four dimensions the pointwise kinetic density obeys the exact identity
\begin{equation}
 \frac12\sqrt{|g|}\,g^{\mu\nu}
 \langle D_\mu\Theta,D_\nu\Theta\rangle_\sharp
 =2\mathcal N_0\sqrt{|g|}.
 \label{eq:volume-collapse}
\end{equation}
Thus the fully metric-locked kinetic scalar is a volume term.  Its composite
second variation must be computed as the Hessian of the volume functional
with the implicit $D$-composite constraints; it is not obtained by silently
reusing the fixed-background Hessian of Theorem~\ref{thm:fixed}.
\end{proposition}

\begin{proof}
The metric lock gives
$\langle D_\mu\Theta,D_\nu\Theta\rangle_\sharp
=\mathcal N_0g_{\mu\nu}$.  Contracting with $g^{\mu\nu}$ gives
$4\mathcal N_0$, and the prefactor one half yields
Eq.~\eqref{eq:volume-collapse}.
\end{proof}

\section{Why the six-dimensional resonance is not the principal bundle}
The frozen $D$-composite analysis gives a $16\times16$ matrix
$A(s,\lambda)$ linear in $s$, with
\begin{equation}
 A^3=qA^2,\qquad q=\lambda\cdot s,\qquad
 \det(I-A)=(1-q)^6.
 \label{eq:dcomp}
\end{equation}
At $q=1$, $\ker(I-A)$ is six-dimensional and has an injective curl map.

\begin{theorem}[Nonconic resonance no-go]
\label{thm:nonconic}
The singular set $q=1$ of the full mixed-order symbol $I-A$ is not conic in
covector space.  In particular, if $(s,\lambda)$ obeys $q=1$, then under
$s\mapsto cs$,
\begin{equation}
 \det\bigl(I-A(cs,\lambda)\bigr)=(1-c)^6,
\end{equation}
which is nonzero for generic $c\ne1$.  Therefore the $q=1$ eigenspace cannot,
as it stands, be the characteristic bundle of a homogeneous principal symbol.
The actual first-order principal symbol of the differential system
$I-A(\partial,\lambda)$ is $-A(s,\lambda)$; at generic exact points it has
rank nine rather than full rank sixteen.  It is neither a minimal scalar
Laplace symbol nor an independently established six-component Laplace
operator.
\end{theorem}

\begin{proof}
Linearity gives $A(cs,\lambda)=cA(s,\lambda)$.  Substitution in the exact
determinant identity yields the displayed radial dependence.  Characteristic
sets of homogeneous principal symbols are invariant under nonzero radial
rescaling, whereas $q=1$ is not.  The exact generic rank-nine statement is
verified directly from the same corrected $A(s,\lambda)$ used by the existing
$D$-composite audit.
\end{proof}

\begin{corollary}[Meaning of the $q=1$ sector]
The six modes are a finite-scale real-exponential resonance of the frozen
self-consistency equation.  They remain a potentially important clue to the
connection/anholonomy sector, but they cannot be inserted directly as
``six physical fields'' in the induced-gravity trace.  Their norm, statistics,
gauge quotient and quadratic action must first be derived from the true
Hessian.
\end{corollary}

\section{The exact remaining composite calculation}
Let
\begin{equation}
 \Fcal[\Theta]=\bigl(E[\Theta],g[\Theta],\Omega[\Theta],\ldots\bigr)
\end{equation}
denote every composite object substituted into a candidate single-$\Theta$
action.  Schematically, the second variation of
$S(\Theta,\Fcal[\Theta])$ contains
\begin{align}
 \Hess_{\rm comp}
 ={}&S_{\Theta\Theta}
 +S_{\Theta\Fcal}\Fcal_\Theta
 +\Fcal_\Theta^{\!*}S_{\Fcal\Theta}
 +\Fcal_\Theta^{\!*}S_{\Fcal\Fcal}\Fcal_\Theta
 +S_{\Fcal}\Fcal_{\Theta\Theta}.
 \label{eq:chain}
\end{align}
Equation~\eqref{eq:chain} is schematic operator notation, but it records the
terms that a fixed-background variation omits.  Since $E$ already contains a
first derivative of $\Theta$ and a composite Levi--Civita connection contains
derivatives of $E$, the differential order and degeneracy of
$\Hess_{\rm comp}$ cannot be inferred from Theorem~\ref{thm:fixed}.

The remaining decision procedure is now finite and explicit:
\begin{enumerate}
 \item choose one candidate single-$\Theta$ action without an imported
 Hilbert--Palatini term;
 \item compute Eq.~\eqref{eq:chain} at a declared vacuum;
 \item identify diffeomorphism, Lorentz and any internal zero modes;
 \item add a declared gauge fixing and its ghost operator;
 \item eliminate only algebraic auxiliary modes;
 \item determine the differential order and the full principal symbol on the
 physical quotient;
 \item if it is second-order and scalar Laplace type after Euclidean
 continuation, compute $\mathcal E$, the supertrace mode count and the
 $a_2$ coefficient.
\end{enumerate}
A failure at step 6 is a genuine architecture result, not a failed heat-kernel
calculation: it would force a nonminimal/pseudodifferential treatment or a
revision such as the explicitly separated profile-field candidate.

\begin{tcolorbox}[colback=green!6!white,colframe=green!55!black,
title=Machine-verified B-tier status]
\textbf{GAP-10D-HESS-FIXED: CLOSED [L1].}
The fixed-background quadratic $\Theta$ kinetic operator has the scalar
principal symbol Eq.~\eqref{eq:scalar-symbol}; its Euclideanized operator is
of minimal Laplace type under the assumptions stated above.

\medskip
\textbf{GAP-10D-RES-PRINCIPAL: CLOSED AS NO-GO [L1].}
The $q=1$ six-dimensional $D$-composite resonance is not conic and cannot be
identified directly with the principal-symbol bundle used in a Gilkey trace.

\medskip
\textbf{GAP-10D-HESS-COMP: NARROWED, not closed.}
The exact remaining lemma is the complete gauge-fixed second variation
Eq.~\eqref{eq:chain}, followed by the physical symbol, mode and ghost count.
This note does not derive the Einstein coefficient, Newton's constant, or an
unconditional GR equation from the canonical single-$\Theta$ action.

\medskip
These diagnostic subgap labels are B-tier machine-verified results.  They do
not alter the signed A-tier claims ledger until the author performs substantive
review and explicitly promotes the wording.
\end{tcolorbox}

\begin{thebibliography}{9}
\bibitem{vassilevich2003}
D.~V.~Vassilevich, ``Heat kernel expansion: user's manual,''
\emph{Phys. Rept.} \textbf{388}, 279--360 (2003).
\bibitem{gilkey1995}
P.~B.~Gilkey, \emph{Invariance Theory, the Heat Equation and the
Atiyah--Singer Index Theorem}, 2nd ed., CRC Press (1995).
\end{thebibliography}
\end{document}
